% Section 4.5: Breaking the Stalemate
% Transition section: motivates the globular cluster analysis (Paper 3)

\section{Breaking the Stalemate: Independent Constraints from Globular Clusters}
\label{sec:4.5}

The analyses discussed in the preceding sections share a common vulnerability: they all operate directly within the complex gamma-ray environment of the Galactic Center, where the results depend sensitively on the adopted interstellar emission model, masking procedure, and treatment of the Fermi Bubbles. Whether one uses template fitting, non-Poissonian statistics, or neural network-based inference, the systematic uncertainties associated with the Galactic diffuse emission propagate into every conclusion about the nature of the GCE.

To make further progress, independent constraints from environments that do not suffer from these systematics are needed.

Globular clusters are well suited for this purpose. The high stellar densities in these systems drive frequent dynamical interactions, producing large populations of low-mass X-ray binaries and, through recycling, millisecond pulsars~\cite{Amerio:2024qor}. The number of MSPs in a given cluster correlates with observable macroscopic properties --- the stellar encounter rate $\Gamma_e$ and the visible-band luminosity $L_V$ (a proxy for stellar mass) --- allowing the total MSP content to be estimated from quantities that are measured independently of the gamma-ray data~\cite{Amerio:2024qor}. Because globular clusters are discrete, spatially isolated systems far from the confused Galactic plane, their gamma-ray emission can be measured cleanly, without the severe interstellar emission model uncertainties that affect analyses of the GCE.

The age of the stellar populations in globular clusters is especially relevant to the GCE debate. The MSP hypothesis requires the pulsars responsible for the excess to be systematically fainter than those observed in the Galactic Plane, since otherwise Fermi-LAT should have already resolved many of them individually~\cite{Hooper:2016rap, Holst:2024fvb}. One natural explanation is that the Galactic Bulge formed its stars early, so its pulsars are very old and have lost much of their rotational kinetic energy, becoming intrinsically dimmer~\cite{Amerio:2024qor}. Globular clusters test this hypothesis directly: their stellar populations are comparably old to the Bulge, and if extreme age naturally produces faint MSPs, that faintness should be observable in these systems as well.

Building on an earlier study by Hooper and Linden~\cite{Hooper:2016rap} that used 7 years of data, in Amerio, Hooper, and Linden~\cite{Amerio:2024qor} we measured the gamma-ray luminosity function of MSPs across the Milky Way's globular cluster system using 15.8 years of Fermi-LAT data. We detected statistically significant gamma-ray emission from 56 of the 157 known globular clusters, 8 of which had not appeared in previous source catalogs. By fitting the observed gamma-ray luminosities of 87 clusters against predictions based on their stellar encounter rates and visible luminosities, we constrained the underlying log-normal luminosity function. The analysis consistently favors a mean pulsar luminosity of $\langle L_\gamma \rangle \sim (1{-}8) \times 10^{33}$~erg/s (integrated between 0.1 and 100~GeV), with a log-normal width $\sigma_L \sim 1.4{-}2.8$. These values are fully consistent with the luminosity function measured independently for MSPs in the Galactic Plane~\cite{Holst:2024fvb}, demonstrating that old stellar populations do not produce systematically fainter pulsars.

The implications for the GCE are direct. If the excess is generated by MSPs with the luminosity function measured in globular clusters, Fermi-LAT's source catalogs should contain $N_\mathrm{MSP} \sim 17{-}37$ individually resolved pulsars from the Inner Galaxy population. Only three MSP candidates have been identified in this region~\cite{Holst:2024fvb}, and systematic radio follow-up campaigns targeting unassociated Fermi sources have not uncovered a hidden population~\cite{Amerio:2024qor}. The GCE can therefore originate from MSPs only if they are significantly less luminous, on average, than those found in any other observed environment.

This independent constraint connects directly to the recent findings of List et al.~\cite{List:2025qbx}. As discussed in Section~\ref{sec:4.4.3}, their energy-dependent CNN analysis pushes the source-count distribution of any putative GCE point sources to extremely faint levels, requiring on the order of $2 \times 10^5$ sources --- each producing far less than one detected photon --- to account for the excess. Such an ultra-faint population would naturally explain the absence of individually resolved pulsars. However, the globular cluster results remove the astrophysical motivation for its existence: if old age does not produce the required ultra-faint MSP populations, no known mechanism can generate $\mathcal{O}(10^5)$ sources of the necessary faintness in the Inner Galaxy. The tension identified from globular clusters and the tension from the energy-dependent SCD analysis thus converge from independent lines of evidence.

The remainder of this chapter presents our analysis in full. The following sections reproduce the published paper~\cite{Amerio:2024qor}, beginning from the data selection and methodology.
